% PART II: Сбор данных (Stages 1-2)
% Source: docs/03-collectors.md, docs/01-data-sources.md, docs/04-analysts.md

\section{NewsScout: агрегация сигналов}

Первый этап сбора информации в пайплайне Delphi Press основан на агенте~\texttt{NewsScout}, чья задача~--- собрать многоканальные сигналы из различных источников и сформировать унифицированный набор записей для последующей обработки.

% Source: 03-collectors.md:477-489

\subsection{Архитектура сбора}

\texttt{NewsScout} реализует двухисточниковую стратегию:

\begin{enumerate}
    \item \textbf{RSS-агрегация}: параллельная загрузка 20--30 RSS-фидов из мировых, региональных и тематических источников
    \begin{itemize}
        \item Глобальные источники (Reuters, AP, BBC, Al Jazeera и др.) --- 10--15 фидов
        \item Региональные (ТАСС, РИА для русскоязычных; CNN, NYT для англоязычных) --- 5--10
        \item Тематические (финансы, политика, наука) --- 5--10
    \end{itemize}

    \item \textbf{Веб-поиск}: параллельное выполнение 3--5 поисковых запросов через API Exa/Jina MCP
    \begin{itemize}
        \item Общий запрос: \emph{major news events {target\_date} world}
        \item Региональный: \emph{news {region\_of\_outlet} {target\_date}}
        \item Тематический: \emph{scheduled events politics economy {target\_date}}
        \item Outlet-specific: \emph{topics covered by {outlet} this week}
        \item Трендовый: \emph{breaking news today {current\_date}} (для контекста)
    \end{itemize}
\end{enumerate}

Оба источника обрабатываются параллельно с помощью \texttt{asyncio.gather()}. Если один источник недоступен (например, RSS-фид не отвечает в течение 30 секунд), пайплайн продолжает работу с другим, логируя ошибку. Критическое условие: если оба источника полностью недоступны, возвращается исключение.

% Source: 03-collectors.md:542-601

\subsection{Дедупликация и релевантность}

После объединения результатов из обоих источников выполняется дедупликация по URL:

\begin{itemize}
    \item Для каждого сигнала вычисляется нормализованный URL (строчные буквы, удаление trailing slash)
    \item При обнаружении дубля сохраняется сигнал с более высоким \texttt{relevance\_score}
\end{itemize}

Релевантность сигнала рассчитывается на основе:
\begin{itemize}
    \item Для RSS-сигналов: свежесть (дата публикации) и надёжность источника
    \item Для веб-поиска: \texttt{score} из API поисковика (нормализуется в диапазон 0--1)
\end{itemize}

Опциональная LLM-классификация активируется для сигналов без категорий. Используется модель \texttt{openai/gpt-4o-mini} с пакетной обработкой (batch по 20 сигналов). Промпт требует классифицировать заголовки по доменам (политика, экономика, наука и т.д.) и извлечь именованные сущности (названия людей, организаций, мест).

% Source: 03-collectors.md:53-129

\subsection{Схема SignalRecord}

Выход \texttt{NewsScout} --- список объектов \texttt{SignalRecord} (100--200 записей), определённых в \texttt{src/schemas/events.py}. Каждый сигнал содержит:

\begin{table}[h!]
\centering
\caption{Поля SignalRecord}
\label{tab:signal_record}
\small
\begin{tabular}{|l|l|l|}
\hline
\textbf{Поле} & \textbf{Тип} & \textbf{Назначение} \\
\hline
\texttt{id} & str & Уникальный ID: \texttt{sig\_\{hash8\}} \\
\texttt{title} & str & Заголовок (макс. 500 символов) \\
\texttt{summary} & str & Краткое содержание (макс. 1000 символов) \\
\texttt{url} & str & URL источника \\
\texttt{source\_name} & str & Название источника (Reuters, ТАСС, BBC и т.д.) \\
\texttt{source\_type} & Enum & RSS, WEB\_SEARCH, SOCIAL, WIRE \\
\texttt{published\_at} & datetime | None & Дата/время публикации (UTC) \\
\texttt{language} & str & ISO 639-1 код (ru, en, zh и т.д.) \\
\texttt{categories} & list[str] & Теги/категории (politics, economy, science и т.д.) \\
\texttt{entities} & list[str] & Именованные сущности (Trump, ЕЦБ, НАТО) \\
\texttt{relevance\_score} & float & Оценка релевантности (0.0--1.0) \\
\hline
\end{tabular}
\end{table}

% Source: 03-collectors.md:744-756

ID сигнала генерируется как SHA-256 хеш конкатенации URL и заголовка, обрезанный до первых 8 символов. Это обеспечивает детерминированность: один и тот же сигнал, встреченный дважды, получит одинаковый ID.

% Source: 03-collectors.md:533-536, 594-599

Завершающий этап --- сортировка по \texttt{relevance\_score} (убывающий порядок) и ограничение на 200 записей. Сигналы отсортированы от наиболее релевантных к наименее релевантным.

---

\section{EventCalendar: структурированные события}

Второй источник информации в Stage 1 --- агент \texttt{EventCalendar}, который выявляет запланированные события на целевую дату пруи с использованием веб-поиска и структурирования через LLM.

% Source: 03-collectors.md:787-830

\subsection{Логика поиска событий}

Процесс разбивается на пять этапов:

\begin{enumerate}
    \item \textbf{Формирование поисковых запросов} (5--8 запросов) по типам событий:
    \begin{itemize}
        \item Политика: \emph{political events scheduled {date}}, \emph{parliamentary sessions votes {date}}
        \item Экономика: \emph{economic data releases {date}}, \emph{central bank meetings {date}}, \emph{earnings reports {date}}
        \item Дипломатия: \emph{international summits meetings {date}}, \emph{UN sessions {date}}
        \item Суды: \emph{major court hearings verdicts {date}}
        \item Культура/спорт: \emph{major events conferences {date}}
    \end{itemize}
    Для русскоязычных СМИ часть запросов дублируется на русском.

    \item \textbf{Параллельный веб-поиск}: выполнение всех запросов асинхронно через Exa/Jina API

    \item \textbf{LLM-структурирование} (модель \texttt{openai/gpt-4o-mini}): извлечение из сырых результатов структурированных событий с полями title, description, event\_type, certainty, location, participants, potential\_impact

    \item \textbf{Дедупликация}: fuzzy matching по заголовкам (ratio > 0.8 по Levenshtein), совпадение типа и даты события

    \item \textbf{Оценка новостной значимости} (модель \texttt{anthropic/claude-sonnet-4}): ранжирование по \texttt{newsworthiness} с учётом вероятности освещения в выбранном издании и исключительности события
\end{enumerate}

% Source: 03-collectors.md:148-228

\subsection{Типология и уверенность}

Типы событий охватывают основные домены публичной жизни:

\begin{equation}
\text{EventType} \in \{\text{POLITICAL, ECONOMIC, DIPLOMATIC, JUDICIAL, MILITARY, CULTURAL, SCIENTIFIC, SPORTS, OTHER}\}
\end{equation}

Каждое событие сопровождается оценкой уверенности, которая отражает степень надёжности информации:

\begin{table}[h!]
\centering
\caption{Уровни уверенности в наступлении события}
\label{tab:event_certainty}
\begin{tabular}{|l|l|}
\hline
\textbf{Уровень} & \textbf{Интерпретация} \\
\hline
CONFIRMED & Официально подтверждено (расписание на сайте, пресс-релиз) \\
LIKELY & Высокая вероятность (по расписанию, анонсировано) \\
POSSIBLE & Возможно (слухи, предварительные анонсы) \\
SPECULATIVE & Спекулятивно (восстановлено из паттернов: ежемесячные отчёты) \\
\hline
\end{tabular}
\end{table}

% Source: 03-collectors.md:156-228

\subsection{Схема ScheduledEvent}

Выход \texttt{EventCalendar} --- список объектов \texttt{ScheduledEvent} (3--10 записей на типичный день, до 50 в день повышенной активности):

\begin{table}[h!]
\centering
\caption{Поля ScheduledEvent}
\label{tab:scheduled_event}
\small
\begin{tabular}{|l|l|l|}
\hline
\textbf{Поле} & \textbf{Тип} & \textbf{Назначение} \\
\hline
\texttt{id} & str & Уникальный ID: \texttt{evt\_\{hash\}} \\
\texttt{title} & str & Название события (макс. 300 символов) \\
\texttt{description} & str & Подробное описание (макс. 500 символов) \\
\texttt{event\_date} & date & Дата события (совпадает с target\_date) \\
\texttt{event\_type} & EventType & Тип события (political, economic и т.д.) \\
\texttt{certainty} & EventCertainty & Уровень уверенности (confirmed/likely/possible/speculative) \\
\texttt{location} & str & Место проведения \\
\texttt{participants} & list[str] & Ключевые участники/организации \\
\texttt{potential\_impact} & str & Потенциальное влияние на новостную повестку \\
\texttt{source\_url} & str & URL источника информации \\
\texttt{newsworthiness} & float & Оценка новостной значимости (0.0--1.0) \\
\hline
\end{tabular}
\end{table}

Значение \texttt{newsworthiness} = 1.0 означает, что событие гарантированно попадёт в топе новостей выбранного издания.

---

\section{OutletHistorian: профиль издания}

Третий источник данных в Stage 1 --- агент \texttt{OutletHistorian}, который анализирует стилевые характеристики и редакционную позицию целевого СМИ на основе исторических публикаций.

% Source: 03-collectors.md:1021-1076

\subsection{Методология анализа}

Анализ разбивается на три параллельных LLM-вызова:

\begin{enumerate}
    \item \textbf{Анализ стиля заголовков} (HeadlineStyle) --- 30 последних заголовков издания
    \item \textbf{Анализ стиля текстов} (WritingStyle) --- 10 первых абзацев статей
    \item \textbf{Анализ редакционной позиции} (EditorialPosition) --- полный контекст 50--100 статей за 30 дней
\end{enumerate}

Все три вызова выполняются одновременно через \texttt{asyncio.gather()}, минимизируя общее время выполнения. Модель используется единая: \texttt{anthropic/claude-sonnet-4} (для высокой точности лингвистического анализа).

Кеширование: если профиль для издания был создан менее 7 дней назад, он возвращается из кеша без переанализа. Это обосновано тем, что стилевые характеристики СМИ меняются медленно.

% Source: 03-collectors.md:323-383

\subsection{HeadlineStyle: метрики заголовков}

Характеристики стиля заголовков с точки зрения формы и содержания:

\begin{table}[h!]
\centering
\caption{Показатели HeadlineStyle}
\label{tab:headline_style}
\footnotesize
\begin{tabular}{|l|l|l|}
\hline
\textbf{Метрика} & \textbf{Тип} & \textbf{Описание} \\
\hline
\texttt{avg\_length\_chars} & int & Средняя длина заголовка (символы) \\
\texttt{avg\_length\_words} & int & Средняя длина заголовка (слова) \\
\texttt{uses\_colons} & bool & Двоеточие для разделения темы и деталя \\
\texttt{uses\_quotes} & bool & Цитаты в заголовках \\
\texttt{uses\_questions} & bool & Вопросительные заголовки \\
\texttt{uses\_numbers} & bool & Частое использование цифр (5 причин...) \\
\texttt{capitalization} & str & Стиль: sentence\_case, title\_case, all\_caps\_first\_word, lowercase \\
\texttt{vocabulary\_register} & str & Регистр: formal, neutral, colloquial, technical, mixed \\
\texttt{emotional\_tone} & str & Тон: neutral, alarming, optimistic, dramatic, ironic, dry \\
\texttt{common\_patterns} & list[str] & 3+ повторяющихся паттерна (например, \{Name\}: statement...) \\
\hline
\end{tabular}
\end{table}

% Source: 03-collectors.md:386-422

\subsection{WritingStyle: характеристики текстов}

Стилевые характеристики основного содержания статей:

\begin{table}[h!]
\centering
\caption{Показатели WritingStyle}
\label{tab:writing_style}
\footnotesize
\begin{tabular}{|l|l|l|}
\hline
\textbf{Поле} & \textbf{Тип} & \textbf{Описание} \\
\hline
\texttt{first\_paragraph\_style} & str & inverted\_pyramid, narrative, analytical, quote\_lead \\
\texttt{avg\_first\_paragraph\_sentences} & int & Среднее число предложений в первом абзаце \\
\texttt{avg\_first\_paragraph\_words} & int & Среднее число слов \\
\texttt{attribution\_style} & str & source\_first, source\_last, inline \\
\texttt{uses\_dateline} & bool & Наличие датлайна (МОСКВА, 2 апреля ---) \\
\texttt{paragraph\_length} & str & short, medium, long \\
\hline
\end{tabular}
\end{table}

\textit{Inverted pyramid} (перевёрнутая пирамида) --- классический журналистский стиль, когда самая важная информация (кто-что-где-когда) размещена в первом абзаце.

% Source: 03-collectors.md:425-472

\subsection{EditorialPosition: редакционная позиция}

Редакционная позиция и предпочтения издания определяют, какие темы получают приоритет и как они подаются:

\begin{table}[h!]
\centering
\caption{Компоненты EditorialPosition}
\label{tab:editorial_position}
\footnotesize
\begin{tabular}{|l|l|l|}
\hline
\textbf{Поле} & \textbf{Тип} & \textbf{Содержание} \\
\hline
\texttt{tone} & ToneProfile & neutral, conservative, liberal, sensationalist, analytical, official, oppositional \\
\texttt{focus\_topics} & list[str] & Приоритетные темы (внутренняя политика, экономика РФ и т.д.) \\
\texttt{avoided\_topics} & list[str] & Обходимые стороной темы \\
\texttt{framing\_tendencies} & list[str] & Типичные фреймы: pro\_government, market\_oriented, human\_interest, conflict\_frame \\
\texttt{source\_preferences} & list[str] & Предпочтительные источники (официальные лица, МИД, анонимные) \\
\texttt{stance\_on\_current\_topics} & dict & Позиция по ключевым темам (тема: позиция) \\
\texttt{omissions} & list[str] & Систематически пропускаемые темы \\
\hline
\end{tabular}
\end{table}

% Source: 03-collectors.md:289-320

\subsection{OutletProfile: итоговая структура}

Полный профиль издания объединяет все три анализа:

\begin{table}[h!]
\centering
\caption{Структура OutletProfile}
\label{tab:outlet_profile}
\footnotesize
\begin{tabular}{|l|l|l|}
\hline
\textbf{Поле} & \textbf{Тип} & \textbf{Назначение} \\
\hline
\texttt{outlet\_name} & str & Каноническое название издания \\
\texttt{outlet\_url} & str & Основной URL издания \\
\texttt{language} & str & ISO 639-1 код (ru, en и т.д.) \\
\texttt{headline\_style} & HeadlineStyle & Объект с метриками заголовков \\
\texttt{writing\_style} & WritingStyle & Объект с метриками текстов \\
\texttt{editorial\_position} & EditorialPosition & Объект с редакционной позицией \\
\texttt{sample\_headlines} & list[str] & 10--30 последних заголовков \\
\texttt{sample\_first\_paragraphs} & list[str] & 5--10 примеров первых абзацев \\
\texttt{analysis\_period\_days} & int & 30 (период анализа) \\
\texttt{articles\_analyzed} & int & Количество проанализированных статей \\
\texttt{analyzed\_at} & datetime & Timestamp проведения анализа \\
\hline
\end{tabular}
\end{table}

Профиль кешируется в SQLite таблице \texttt{outlet\_profiles} с ключом \texttt{outlet\_name} (нормализованное название) и TTL 7 дней.

---

\section{ForesightCollector: рынки предсказаний}

Четвёртый источник информации (дополнительный, на Stage 1.5) --- агент \texttt{ForesightCollector}, собирающий данные из рынков предсказаний для обогащения контекста прогноза.

% Основано на инженерной реализации и документации

\subsection{Polymarket CLOB API}

Основной источник данных --- \emph{Polymarket} (децентрализованная платформа предсказывающих рынков, построенная на блокчейне). Доступ осуществляется через HTTP REST API без требования аутентификации.

\textbf{Ключевые параметры API}:
\begin{itemize}
    \item URL: \texttt{https://clob.polymarket.com}
    \item Endpoints: \texttt{/markets}, \texttt{/prices}, \texttt{/trades}
    \item Rate limit: ~200 запросов/10 секунд
    \item Timeout: 10 секунд
\end{itemize}

Присоединение к рынкам выполняется через \texttt{conditionId} --- уникальный идентификатор условия события на блокчейне. Например, рынок «Will Trump win 2024 election?» может иметь conditionId = \texttt{0x1234...}. Этот ID используется как primary join key при агрегировании данных.

% Основано на Phase 6 completed

\subsection{Сбираемые метрики рынка}

Для каждого релевантного рынка собираются следующие метрики:

\begin{table}[h!]
\centering
\caption{Метрики Polymarket CLOB API}
\label{tab:polymarket_metrics}
\small
\begin{tabular}{|l|l|l|}
\hline
\textbf{Метрика} & \textbf{Источник} & \textbf{Назначение} \\
\hline
\texttt{conditionId} & markets endpoint & Уникальный идентификатор \\
\texttt{question} & markets & Формулировка предсказания \\
\texttt{volume} & prices & Объём торговли (всего) \\
\texttt{open\_interest} & prices & Открытый интерес (активные позиции) \\
\texttt{bid\_price} & prices & Лучшая цена покупки (YES) \\
\texttt{ask\_price} & prices & Лучшая цена продажи (YES) \\
\texttt{mid\_price} & computed & (bid + ask) / 2 \\
\texttt{spread} & computed & ask - bid (индикатор ликвидности) \\
\texttt{resolution\_date} & markets & Дата закрытия рынка \\
\texttt{last\_updated} & prices & Timestamp последнего обновления \\
\hline
\end{tabular}
\end{table}

\textit{Spread} (разница между ценами покупки и продажи) служит индикатором ликвидности рынка: узкий spread указывает на активную торговлю и уверенность участников.

% Основано на инженерной реализации

\subsection{Graceful degradation}

\textbf{Metaculus API} (альтернативный источник) в настоящее время недоступен из-за требования BENCHMARKING-уровня доступа (статус: запрошено 2026-03-29). Код пайплайна готов к интеграции, но на практике возвращает пустой список при попытке доступа.

Если Polymarket API недоступен (HTTP 5xx, timeout, rate limit):
\begin{enumerate}
    \item Retry с exponential backoff (3 попытки, макс. 30 сек)
    \item При повторных ошибках: логирование warning и возврат пустого списка `[]`
    \item Pipeline продолжает работу без данных рынков (деградация качества, но не критично)
\end{enumerate}

Успешный API-вызов возвращает список рынков с метриками. Релевантность рынка определяется совпадением между question/description рынка и событиями из EventCalendar (косинусное сходство эмбеддингов или простое ключевое слово matching).

---

\section{Идентификация событий: HDBSCAN-кластеризация}

Stage 2 пайплайна агломерирует 100--200 сигналов и 3--10 запланированных событий в 20 структурированных событийных нитей через плотностную кластеризацию.

% Source: 04-analysts.md:785-836

\subsection{Логика процесса}

Процесс разбивается на следующие этапы:

\begin{enumerate}
    \item \textbf{Подготовка текстов}: конкатенация title + summary для каждого SignalRecord

    \item \textbf{Эмбеддинг}: batch-эмбеддинг через LLM API (OpenAI embeddings или Voyage AI). Результат: матрица размером (N, 1536), где N = количество сигналов

    \item \textbf{Кластеризация HDBSCAN}: плотностная кластеризация эмбеддингов в пространстве косинусного сходства

    \item \textbf{LLM-лейблинг}: для каждого кластера генерируется название и summary через LLM

    \item \textbf{Интеграция ScheduledEvent}: запланированные события привязываются к кластерам на основе семантического сходства

    \item \textbf{Scoring}: расчёт \texttt{significance\_score} по многофакторной формуле

    \item \textbf{Ранжирование}: отбор top-20 по significance\_score

    \item \textbf{Trajectory Analysis}: для каждой нити анализируется текущее состояние, моментум, три сценария развития

    \item \textbf{Cross-Impact Matrix}: оценка перекрёстных влияний между событийными нитями
\end{enumerate}

% Source: 04-analysts.md:804-818

\subsection{Гиперпараметры HDBSCAN}

\begin{table}[h!]
\centering
\caption{Гиперпараметры кластеризации}
\label{tab:hdbscan_params}
\begin{tabular}{|l|l|l|}
\hline
\textbf{Параметр} & \textbf{Значение} & \textbf{Интерпретация} \\
\hline
\texttt{min\_cluster\_size} & 3 & Минимальный размер кластера (события, упомянутые в 3+ сигналах) \\
\texttt{min\_samples} & 2 & Минимальный размер neighbourhood для плотности \\
\texttt{metric} & cosine & Метрика расстояния в пространстве эмбеддингов \\
\hline
\end{tabular}
\end{table}

\texttt{min\_cluster\_size} = 3 означает, что кластер формируется только если он содержит минимум 3 сигнала с близкими эмбеддингами (косинусное сходство). Это фильтрует шум (случайные совпадения одного-двух заголовков) и выделяет только события, получившие достаточное освещение в множественных источниках.

% Source: 04-analysts.md:812-818

\subsection{Формула значимости события}

Значимость события рассчитывается как взвешенная сумма пяти компонент:

\begin{equation}
\text{significance\_score} = 0.30 \times \text{importance} + 0.25 \times \text{cluster\_size\_norm} + 0.20 \times \text{recency} + 0.15 \times \text{source\_diversity} + 0.10 \times \text{entity\_prominence}
\label{eq:significance}
\end{equation}

Компоненты определяются следующим образом:

\begin{enumerate}
    \item \textbf{importance} (вес 0.30) --- оценка LLM медийной важности события в контексте выбранного издания (0.0--1.0). Высокое значение для событий, которые гарантированно будут в топе новостей.

    \item \textbf{cluster\_size\_norm} (вес 0.25) --- нормализованный размер кластера. Если кластер содержит $k$ сигналов, а максимальный размер $k_{\max}$, то этот компонент = $k / k_{\max}$ (0.0--1.0). Большие кластеры указывают на высокую новостную активность.

    \item \textbf{recency} (вес 0.20) --- свежесть сигналов в кластере. Сигналы, опубликованные в последние 24 часа, получают полный балл 1.0. Более старые сигналы штрафуются экспоненциально: $\exp(-t / \tau)$, где $t$ --- возраст (дни), $\tau = 3$ (временная постоянная).

    \item \textbf{source\_diversity} (вес 0.15) --- количество уникальных источников в кластере, нормализованное: $\min(1.0, |\text{sources}| / 5)$. Событие, покрытое 5+ различными источниками, получает максимальный балл.

    \item \textbf{entity\_prominence} (вес 0.10) --- число упоминаний высокопрофильных сущностей (главы государств, ЕЦБ, ООН и т.д.) в сигналах кластера, нормализованное к [0, 1]. Для типичного кластера этот компонент невелик, но может быть решающим для крупных политических событий.
\end{enumerate}

% Source: 04-analysts.md:769-836

\subsection{Выход: EventThread[]}

Результат кластеризации --- список объектов \texttt{EventThread} (до 20 записей), упорядоченный по \texttt{significance\_score}:

\begin{table}[h!]
\centering
\caption{Поля EventThread}
\label{tab:event_thread}
\footnotesize
\begin{tabular}{|l|l|l|}
\hline
\textbf{Поле} & \textbf{Тип} & \textbf{Содержание} \\
\hline
\texttt{id} & str & Уникальный идентификатор нити \\
\texttt{title} & str & Название события (сгенерировано LLM из кластера) \\
\texttt{summary} & str & Summary событийной нити \\
\texttt{signal\_ids} & list[str] & ID входящих сигналов \\
\texttt{scheduled\_events} & list[str] & ID привязанных запланированных событий \\
\texttt{cluster\_size} & int & Количество сигналов в кластере \\
\texttt{significance\_score} & float & Итоговая значимость (0.0--1.0) \\
\texttt{sources} & list[str] & Уникальные источники сигналов \\
\texttt{entities} & list[str] & Агрегированные именованные сущности \\
\texttt{published\_at} & datetime & Время публикации самого старого сигнала \\
\texttt{last\_updated} & datetime & Время самого свежего сигнала \\
\texttt{trajectory} & EventTrajectory | None & Анализ траектории (Stage 3) \\
\hline
\end{tabular}
\end{table}

Каждая EventThread представляет связанный набор новостей об одном событии или процессе, готовый для анализа на последующих стадиях пайплайна (Stage 3: тренд-анализ, Stage 4--5: Дельфи-прогнозирование).

---

\subsection{Обработка ошибок и граничные случаи}

\begin{table}[h!]
\centering
\caption{Стратегии обработки ошибок в Stage 2}
\label{tab:stage2_errors}
\footnotesize
\begin{tabular}{|l|l|}
\hline
\textbf{Ситуация} & \textbf{Действие} \\
\hline
< 10 сигналов на входе & Вернуть как есть + warning; HDBSCAN может не сформировать кластеры \\
Эмбеддинг API unavailable & Fallback: использовать простую текстовую дистанцию (Jaccard) \\
0 кластеров найдено & Трактовать каждый сигнал как отдельный EventThread \\
LLM не парсит JSON при лейблинге & Retry 1 раз; использовать generic title при неудаче \\
\hline
\end{tabular}
\end{table}

